%-------------------------
% Chen-Optimized CV for EB2-NIW Petition
% Based on existing CV content; avoids adding unverified claims
%------------------------

\documentclass[letterpaper,11pt]{article}

\usepackage{latexsym}
\usepackage[empty]{fullpage}
\usepackage{titlesec}
\usepackage{marvosym}
\usepackage[usenames,dvipsnames]{color}
\usepackage{verbatim}
\usepackage{enumitem}
\usepackage[hidelinks]{hyperref}
\usepackage{fancyhdr}
\usepackage[english]{babel}
\usepackage{tabularx}
\input{glyphtounicode}

\pagestyle{fancy}
\fancyhf{}
\fancyfoot{}
\renewcommand{\headrulewidth}{0pt}
\renewcommand{\footrulewidth}{0pt}

\setlength{\footskip}{4.08003pt}
\addtolength{\oddsidemargin}{-0.5in}
\addtolength{\evensidemargin}{-0.5in}
\addtolength{\textwidth}{1in}
\addtolength{\topmargin}{-.5in}
\addtolength{\textheight}{1.0in}

\urlstyle{same}
\raggedbottom
\raggedright
\setlength{\tabcolsep}{0in}

\titleformat{\section}{
  \vspace{-4pt}\scshape\raggedright\large
}{}{0em}{}[\color{black}\titlerule \vspace{-5pt}]

\pdfgentounicode=1

\newcommand{\resumeItem}[1]{
  \item\small{{#1 \vspace{-2pt}}}
}

\newcommand{\resumeSubheading}[4]{
  \vspace{-2pt}\item
    \begin{tabular*}{0.97\textwidth}[t]{l@{\extracolsep{\fill}}r}
      \textbf{#1} & #2 \\
      \textit{\small#3} & \textit{\small #4} \\
    \end{tabular*}\vspace{-7pt}
}

\newcommand{\resumeProjectHeading}[2]{
    \item
    \begin{tabular*}{0.97\textwidth}{l@{\extracolsep{\fill}}r}
      \small#1 & #2 \\
    \end{tabular*}\vspace{-7pt}
}

\newcommand{\resumeSubItem}[1]{\resumeItem{#1}\vspace{-4pt}}
\renewcommand\labelitemii{$\vcenter{\hbox{\tiny$\bullet$}}$}
\newcommand{\resumeSubHeadingListStart}{\begin{itemize}[leftmargin=0.15in, label={}]}
\newcommand{\resumeSubHeadingListEnd}{\end{itemize}}
\newcommand{\resumeItemListStart}{\begin{itemize}}
\newcommand{\resumeItemListEnd}{\end{itemize}\vspace{-5pt}}

\begin{document}

\begin{center}
    \textbf{\Huge \scshape Ibrahim Odat} \\ \vspace{2pt}
    \small\textit{Ph.D. Student in Computer Science \& Informatics $|$ Cybersecurity Researcher $|$ AI Security for Autonomous UAV and Multi-Agent Systems} \\ \vspace{3pt}
    \small (248) 237-5426 $|$
    \href{mailto:ibrahimodatt@gmail.com}{\underline{ibrahimodatt@gmail.com}} $|$
    Michigan, USA \\ \vspace{1pt}
    \href{https://3odat.github.io}{\underline{3odat.github.io}} $|$
    \href{https://linkedin.com/in/ibrahim-odat}{\underline{linkedin.com/in/ibrahim-odat}} $|$
    \href{https://scholar.google.com/citations?user=U5HR4EcAAAAJ&hl=en}{\underline{Google Scholar}}
\end{center}

\section{Proposed Endeavor Summary}
\begin{itemize}[leftmargin=0.15in, label={}]
\small{\item{
Ph.D. student in Computer Science and Informatics at Oakland University specializing in the security of AI-enabled autonomous drones (UAVs), LLM-based cyber-physical systems, and multi-agent architectures. My proposed endeavor is to advance the security and resilience of autonomous UAV and multi-agent systems through research on adversarial attacks, secure memory and retrieval mechanisms, and defensive architectures for mission-critical autonomous platforms. My work includes peer-reviewed research on LLM-enabled UAV communication attacks and ongoing research on memory-poisoning threats in LLM-based multi-agent UAV systems.
}}
\end{itemize}

\section{Education}
\resumeSubHeadingListStart
\resumeSubheading
{Oakland University}{Rochester, MI}
{Doctor of Philosophy (Ph.D.), Computer Science \& Informatics $|$ GPA: 3.93/4.0}{Expected Dec. 2028}
\resumeItemListStart
\resumeItem{Research focus: security of LLM-enabled autonomous drones (UAVs) and multi-agent systems}
\resumeItem{Selected coursework: Large Language Models, Cyber Laws \& Digital Forensics, Integrated Computing Systems, Advanced Data Structures \& Algorithms}
\resumeItemListEnd

\resumeSubheading
{Oakland University}{Rochester, MI}
{Master of Science (M.S.), Cybersecurity $|$ GPA: 4.0/4.0}{May 2022 -- Dec. 2023}
\resumeSubHeadingListEnd

\section{Selected Peer-Reviewed Publications}
\resumeSubHeadingListStart
\resumeSubheading
{Net-GPT: A LLM-Empowered Man-in-the-Middle Chatbot for UAV}{2023}
{IEEE/ACM Symposium on Edge Computing (SEC 2023), pp. 287--293}{Peer-reviewed conference publication}
\resumeItemListStart
\resumeItem{Research on the use of large language models to automate man-in-the-middle attack generation against drone-to-ground-station communications}
\resumeItem{The paper has received Google Scholar citations, indicating follow-on interest from researchers in related cyber-physical and communication-system domains}
\resumeItemListEnd

\resumeSubheading
{Heterogeneous Generative Dataset for UASes}{2023}
{IEEE International Conference on Mobility, Operations, Services and Technologies (MOST 2023)}{Peer-reviewed conference publication}
\resumeItemListStart
\resumeItem{Co-authored research on heterogeneous dataset generation for unmanned aerial systems security evaluation}
\resumeItemListEnd
\resumeSubHeadingListEnd

\section{Research Manuscript in Preparation}
\resumeSubHeadingListStart
\resumeSubheading
{Memory as a Control Plane: Poisoning Attacks on LLM Multi-Agent UAV Systems}{2026}
{Research manuscript in preparation with faculty collaborator, Oakland University}{}
\resumeItemListStart
\resumeItem{Investigates how memory poisoning in shared-memory LLM-based UAV systems can influence coordinated mission behavior across sequential missions in simulation}
\resumeItem{Develops and evaluates defensive mechanisms for reducing susceptibility to retrieval- and memory-stage manipulation in multi-agent UAV architectures}
\resumeItem{Introduces a threat model for memory poisoning in LLM-powered multi-agent UAV systems}
\resumeItemListEnd
\resumeSubHeadingListEnd

\section{Research Experience}
\resumeSubHeadingListStart
\resumeSubheading
{Cybersecurity Researcher}{Jan. 2022 -- Present}
{Oakland University}{Rochester, MI}
\resumeItemListStart
\resumeItem{Conducted security research on PX4 Autopilot UAV flight controllers, including analysis of buffer overflows, session hijacking, and MAVLink-related attack surfaces across multiple UAV communication paths}
\resumeItem{Developed AI-assisted threat-detection and anomaly-analysis pipelines using fine-tuned language models, prompt engineering, and local LLM deployment workflows}
\resumeItem{Performed firmware analysis on IoT camera systems, identifying authentication and memory-safety weaknesses and documenting remediation recommendations}
\resumeItem{Simulated adversarial attacks against encryption and intrusion-detection configurations to support improved detection logic and incident response preparation}
\resumeItem{Designed and implemented AeroMind, a multi-agent UAV research platform integrating LLMs, persistent memory, and PX4 SITL simulation for autonomous mission execution and security evaluation}
\resumeItemListEnd
\resumeSubHeadingListEnd

\section{Selected Professional Experience}
\resumeSubHeadingListStart
\resumeSubheading
{Cybersecurity Specialist}{Jan. 2021 -- Jan. 2022}
{ZADD}{Amman, Jordan}
\resumeItemListStart
\resumeItem{Integrated AI-assisted anomaly detection into an enterprise SaaS security workflow to improve analyst triage efficiency}
\resumeItem{Executed vulnerability assessments across enterprise infrastructure and supported remediation efforts aligned with established security-control frameworks}
\resumeItem{Strengthened Windows Active Directory environments by identifying lateral-movement risks and improving access-control posture}
\resumeItem{Operated and tuned IDS/IPS, firewalls, DLP, and SIEM technologies to maintain visibility across enterprise security operations}
\resumeItemListEnd

\resumeSubheading
{Penetration Tester}{Jan. 2020 -- Jan. 2021}
{Futuretec Security Solutions}{Amman, Jordan}
\resumeItemListStart
\resumeItem{Performed web application and API penetration testing using industry-standard offensive security tools, identifying injection, authentication, and misconfiguration vulnerabilities}
\resumeItem{Assessed systems against OWASP Top 10 risks and supported remediation for client security and compliance objectives}
\resumeItem{Developed Python-based automation to support vulnerability discovery and streamline portions of manual testing workflows}
\resumeItem{Prepared remediation reports for technical and management stakeholders and collaborated with development teams on secure deployment practices}
\resumeItemListEnd
\resumeSubHeadingListEnd

\section{Selected Research Projects and Open-Source Artifacts}
\resumeSubHeadingListStart
\resumeProjectHeading
{\textbf{AeroMind Drone Security Testbed} $|$ \emph{LangGraph, PX4 SITL, FastAPI, Python} $|$ \href{https://github.com/3odat/aeromind-uav-security-testbed}{\underline{Repository}}}{}
\resumeItemListStart
\resumeItem{Built a multi-agent UAV research platform with a supervisor agent, autonomous scout drones, and persistent memory to evaluate memory poisoning, cross-agent propagation, and retrieval-stage defenses}
\resumeItemListEnd

\resumeProjectHeading
{\textbf{Net-GPT Evaluation Toolkit} $|$ \emph{Python, LLM Fine-Tuning, MAVLink, Packet Generation} $|$ \href{https://github.com/3odat/net-gpt-evaluation-toolkit}{\underline{Repository}}}{}
\resumeItemListStart
\resumeItem{Developed reusable scripts and evaluation workflows for fine-tuning language models, generating UAV packets, and reproducing the methodology associated with prior peer-reviewed UAV security research}
\resumeItemListEnd

\resumeProjectHeading
{\textbf{PX4 Autopilot Session Hijacking Exploit} $|$ \emph{PX4 Autopilot, MAVLink, MAVSDK, Python} $|$ \href{https://github.com/3odat/px4-session-hijacking-exploit}{\underline{Repository}}}{}
\resumeItemListStart
\resumeItem{Analyzed MAVLink session management and developed a proof-of-concept exploit demonstrating telemetry-stream hijacking and related remediation considerations}
\resumeItemListEnd
\resumeSubHeadingListEnd

\section{Technical Skills}
\begin{itemize}[leftmargin=0.15in, label={}]
\small{\item{
\textbf{AI/LLM Security}{: Adversarial ML, LLM Fine-Tuning (QLoRA), Prompt Engineering, RAG Poisoning, Memory Poisoning, LangGraph, LangChain, Ollama} \\
\textbf{Autonomous Systems}{: ROS2, PX4 Autopilot, Gazebo, MAVSDK, MAVLink, NVIDIA Jetson Nano, YOLOv8} \\
\textbf{Offensive Security}{: Penetration Testing, Exploit Development, Firmware Analysis, Reverse Engineering, Red Teaming, OSINT} \\
\textbf{Frameworks \& Compliance}{: MITRE ATT\&CK, OWASP Top 10, NIST SP 800-171, NIST AI RMF, PCI DSS, ISO 27001} \\
\textbf{Security Tooling}{: Kali Linux, Metasploit, Burp Suite, Wireshark, Nessus, Nmap, SQLMap, SIEM, IDS/IPS} \\
\textbf{Programming \& Infrastructure}{: Python, Bash, C/C++, SQL, FastAPI, Docker, AWS, Active Directory}
}}
\end{itemize}

\section{Certifications}
\begin{itemize}[leftmargin=0.15in, label={}]
\small{\item{
\begin{itemize}
\item \textbf{Certified Ethical Hacker (CEH)} -- EC-Council, 2023
\item \textbf{eLearnSecurity Junior Penetration Tester (eJPT)} -- INE, 2022
\item \textbf{Fortinet Network Security Expert NSE 1 \& NSE 2} -- Fortinet, 2021
\end{itemize}
}}
\end{itemize}

\end{document}
